\section{Method}
\label{sec:method}

We keep the Scaffold-GS scene representation and rasterizer, and replace each shared attribute decoder by a mixture of experts.
A flat MoE is trained first and then converted into a growing tree, so that additional experts are allocated only where residual difficulty remains high.
Default values of all symbols appear in Table~\ref{tab:hparams}.

\subsection{Preliminaries: Scaffold-GS}
\label{sec:prelim}

Scaffold-GS represents a scene by anchors $\mathcal{A}=\{(\mathbf{x}_i,\mathbf{f}_i)\}_{i=1}^{N}$,
where $\mathbf{x}_i\in\mathbb{R}^{3}$ is a position and $\mathbf{f}_i\in\mathbb{R}^{C}$ is a learnable feature.
For a camera, visible anchors produce a view-conditioned input $\mathbf{u}_i=\phi(\mathbf{f}_i,\mathbf{d}_i,\ldots)$,
with $\mathbf{d}_i$ the viewing direction.
Three attribute networks $F_h$, $h\in\{\alpha,\Sigma,c\}$, map $\mathbf{u}_i$ to opacity, covariance, and color of the associated neural Gaussians.
The predicted image is
\begin{equation}
\hat{I}
=
\mathcal{R}\!\left(
\big\{(\alpha_j,\boldsymbol{\Sigma}_j,\mathbf{c}_j)\big\}_{j}
\right),
\label{eq:render}
\end{equation}
where $\mathcal{R}$ denotes differentiable Gaussian rasterization.
We retain $\mathcal{R}$ and the Scaffold-GS attribute-conversion heads, and replace only $\{F_h\}$.
Throughout, the router of each head consumes the anchor feature $\mathbf{f}_i$ and does not use a GNN.

\subsection{Flat Mixture-of-Experts Decoding}
\label{sec:flat}

Each head $h$ consists of $E$ expert MLPs $\{f_e^{(h)}\}_{e=1}^{E}$ and a linear router $r^{(h)}$.
$E$ is the number of specialists that share the work of one Scaffold-GS MLP;
the router maps the anchor feature to a distribution over those specialists.
Experts share a two-layer architecture with a ReLU hidden unit and the same output activation as Scaffold-GS.
The router produces mixing weights
\begin{equation}
\mathbf{g}_i^{(h)}
=
\operatorname{softmax}\!\big(
r^{(h)}(\mathbf{f}_i)+\boldsymbol{\eta}_i^{(h)}
\big)
\in\Delta^{E-1}.
\label{eq:flat-gates}
\end{equation}
During dense training we add i.i.d.\ Gaussian logit noise $\boldsymbol{\eta}_i^{(h)}\sim\mathcal{N}(\mathbf{0},\sigma_{\eta}^{2}\mathbf{I})$,
where $\sigma_{\eta}$ is a small exploration scale that prevents the router from collapsing onto one expert too early;
otherwise $\boldsymbol{\eta}_i^{(h)}=\mathbf{0}$.

For the first $T_{\mathrm{w}}$ iterations---a dense warm-up---every expert remains active and $\mathbf{w}_i^{(h)}=\mathbf{g}_i^{(h)}$.
After warm-up, Top-$k$ keeps only the $k$ largest gates (the $k$ experts assigned to that anchor) and renormalizes:
\begin{equation}
\begin{aligned}
\mathbf{w}_i^{(h)}
&=
\frac{\operatorname{TopK}(\mathbf{g}_i^{(h)},k)}
{\|\operatorname{TopK}(\mathbf{g}_i^{(h)},k)\|_{1}},
\\
F_h(\mathbf{u}_i^{(h)})
&=
\sum_{e=1}^{E}
w_{i,e}^{(h)}\,f_e^{(h)}(\mathbf{u}_i^{(h)}).
\end{aligned}
\label{eq:topk}
\end{equation}
$\operatorname{TopK}(\cdot,k)$ zeroes all but the $k$ largest entries.
The implementation evaluates every expert and then forms the weighted sum, so Top-$k$ sparsifies the mixture weights rather than skipping expert computation.

\subsection{Growing Tree MoE}
\label{sec:tree}

Each attribute head maintains its own rooted tree: routers sit at internal nodes and experts at leaves $\mathcal{L}_h$.
Stage~B starts from a depth-1 tree with $E$ leaves under the root, copying the flat-MoE experts and root router one-to-one.

At a node $\nu$ with $b_{\nu}$ children, the same dense / Top-$k$ rule as in Eq.~\eqref{eq:topk} is applied with
$k_{\nu}=\min(k,b_{\nu})$.
Top-$k$ is therefore a \emph{node-level} restriction, not a cap on the number of active leaves in the whole tree.
The mixture weight of leaf $\ell$ is the product of mix weights along its path,
\begin{equation}
\begin{aligned}
\pi_i^{(h)}(\ell)
&=
\prod_{(\nu\rightarrow j)\in\operatorname{path}(\ell)}
w_{i,\nu\rightarrow j}^{(h)},
\\
F_h(\mathbf{u}_i^{(h)})
&=
\sum_{\ell\in\mathcal{L}_h}
\pi_i^{(h)}(\ell)\,f_\ell^{(h)}(\mathbf{u}_i^{(h)}).
\end{aligned}
\label{eq:tree-output}
\end{equation}
Dense path weights $\tilde{\pi}_i^{(h)}(\ell)$ are defined analogously from the full softmax gates and are used in auxiliary losses.

A split replaces a leaf $\ell$ by a new linear router and $G$ child experts, where $G$ is the branching factor of a growth event (binary in our experiments).
Child parameters are cloned from the parent and perturbed,
\begin{equation}
\boldsymbol{\theta}_{\ell\rightarrow j}^{(h)}
=
\boldsymbol{\theta}_{\ell}^{(h)}+\boldsymbol{\varepsilon}_{j},
\qquad
\boldsymbol{\varepsilon}_{j}\sim\mathcal{N}(\mathbf{0},\sigma_{\mathrm{split}}^{2}\mathbf{I}),
\label{eq:split}
\end{equation}
so that the $G$ children start as copies of the parent but are not identical;
$\sigma_{\mathrm{split}}$ is the scale of that symmetry-breaking noise.
Each split adds $G-1$ net leaves.
Three hard caps keep trees from growing without bound:
$D_{\max}$ is the maximum depth of any leaf,
$G_{\max}$ is the maximum number of splits allowed on one head,
and $L_{\max}$ is the maximum number of leaves on one head.
At a growth check, at most $H_{\mathrm{evt}}$ heads may split (so opacity, covariance, and color cannot all grow at once), and at most one leaf per selected head is split.

\subsection{Traffic- and Difficulty-Guided Growth}
\label{sec:growth}

Growth uses three training-time statistics---routing traffic, residual reconstruction error, and feature variation---as proxies for where extra capacity may help.
They are not a formal information-gain criterion.
We drop the head index below.

Let $\mathcal{V}_t$ be the visible anchors at iteration $t$, and let $\mathcal{W}$ be the iterations since the last statistics reset of that head.
With mix weights $\pi_i^{(t)}$ (dense or Top-$k$, whichever is used for rendering),
\begin{equation}
\begin{aligned}
m_{\ell}^{(t)}
&=
\sum_{i\in\mathcal{V}_t}\pi_i^{(t)}(\ell),
\qquad
m_{\ell}
=
\sum_{t\in\mathcal{W}}m_{\ell}^{(t)},
\qquad
p_{\ell}
=
\frac{m_{\ell}}{\sum_{\ell'}m_{\ell'}}.
\end{aligned}
\label{eq:traffic}
\end{equation}
Here $m_{\ell}$ is the total routing mass that reached leaf $\ell$ in the window (how much of the scene that leaf is responsible for),
and $p_{\ell}$ is that mass as a fraction of all leaves.
A traffic-weighted average of the \emph{image-level} $\ell_1$ loss $\mathcal{L}_{1}^{(t)}$ is
\begin{equation}
L_{\ell}
=
\frac{\sum_{t\in\mathcal{W}}m_{\ell}^{(t)}\,\mathcal{L}_{1}^{(t)}}{m_{\ell}}.
\label{eq:leaf-loss}
\end{equation}
$L_{\ell}$ is a proxy for how hard the views assigned to $\ell$ are to reconstruct.
Because $\mathcal{L}_{1}^{(t)}$ is shared across leaves of a given view, $L_{\ell}$ ranks leaves through their traffic over $\mathcal{W}$, not through pixel-to-leaf residual assignment.
Feature variation $v_{\ell}$ measures whether the anchors routed to $\ell$ have similar features or mixed ones: it is the traffic-weighted, dimension-averaged variance of $\mathbf{f}_i$,
\begin{equation}
\begin{aligned}
\boldsymbol{\mu}_{\ell}^{(t)}
&=
\frac{\sum_{i\in\mathcal{V}_t}\pi_i^{(t)}(\ell)\,\mathbf{f}_i}{m_{\ell}^{(t)}},
\\
v_{\ell}^{(t)}
&=
\frac{\sum_{i\in\mathcal{V}_t}\pi_i^{(t)}(\ell)\,
\|\mathbf{f}_i-\boldsymbol{\mu}_{\ell}^{(t)}\|_{2}^{2}}{C\,m_{\ell}^{(t)}},
\\
v_{\ell}
&=
\frac{1}{|\mathcal{W}_{\ell}|}\sum_{t\in\mathcal{W}_{\ell}}v_{\ell}^{(t)},
\end{aligned}
\label{eq:leaf-variance}
\end{equation}
where $\mathcal{W}_{\ell}\subseteq\mathcal{W}$ keeps steps with non-negligible mass on $\ell$.

Relative difficulty $R_{\ell}$ is the leaf loss divided by the median leaf loss: values above $1$ mean this leaf is harder than a typical leaf.
The split score multiplies traffic, excess difficulty, and feature spread:
\begin{equation}
\begin{aligned}
R_{\ell}
&=
\frac{L_{\ell}}{\operatorname{median}_{\ell'}L_{\ell'}},
\\
S_{\ell}
&=
p_{\ell}\,
\bigl(1+\max(0,R_{\ell}-1)\bigr)^{\beta}\,
(1+\lambda_{v}v_{\ell}).
\end{aligned}
\label{eq:growth-score}
\end{equation}
Here $\beta\ge 1$ controls how strongly above-median reconstruction error is emphasized,
and $\lambda_{v}\ge 0$ controls how much mixed anchor features raise the score.
A leaf is therefore preferred if it is used often, is harder than the median leaf, and covers diverse features.
Only leaves with recorded traffic enter these statistics.

A structurally eligible leaf (depth and capacity limits) is a candidate only if it also clears four gates
\begin{equation}
p_{\ell}\ge\tau_{p},
\quad
R_{\ell}\ge\tau_{R},
\quad
1+\frac{v_{\ell}}{\operatorname{median}_{\ell'}v_{\ell'}}\ge\tau_{v},
\quad
S_{\ell}\ge\tau_{S}\,
\operatorname{median}_{\ell'\in\mathcal{C}}S_{\ell'},
\label{eq:growth-gates}
\end{equation}
where $\mathcal{C}$ is the set of traffic-supported, structurally eligible leaves of that head.
In words:
$\tau_{p}$ is the minimum traffic share, so rarely used leaves are not split;
$\tau_{R}$ is the minimum relative loss, so easy leaves are not split;
$\tau_{v}$ is the minimum relative feature diversity, so a leaf must mix more varied features than the median;
and $\tau_{S}$ requires the score to exceed the median score by a margin, so only a clearly stronger candidate grows.

To keep the three attribute trees from diverging, a head that currently has fewer leaves than the largest head is treated as \emph{lagging}.
For lagging heads the thresholds $(\tau_{S},\tau_{p},\tau_{R})$ are multiplied by scale factors $\boldsymbol{\gamma}<1$ (a uniform relaxation that makes a smaller tree easier to grow), and the variation gate $\tau_{v}$ is omitted.
An additional relaxation is applied when the color head is the one that lags, since appearance is typically the hardest attribute.
If that head still has no ordinary candidate after another head has grown, a fallback may admit its highest-scoring structurally valid leaf.
Across heads, proposals are ordered by lagging status, then by a fixed head priority (color, opacity, covariance), then by $S_{\ell}$.
The highest-scoring eligible leaf in each of the top $H_{\mathrm{evt}}$ heads is split.

After a successful split, statistics of that head are cleared and newly created parameters are registered with Adam, preserving optimizer state for surviving weights.
Heads that do not split keep their accumulated statistics.
All heads then use dense routing for $T_{\mathrm{d}}$ iterations before resuming Top-$k$,
giving the new children a short window in which they can attract traffic before sparsification.

\begin{figure}[t]
\centering
\includegraphics[width=\linewidth]{Tree_MoE_CVPR_LaTeX/tree_growth_concept.png}
\caption{Tree MoE growth. Each head starts from the flat-MoE experts as root leaves. A selected leaf is replaced by a local router and $G$ child experts. Heads keep separate trees and split budgets; at most $H_{\mathrm{evt}}$ heads split at a given check.}
\label{fig:tree_growth_concept}
\end{figure}

\subsection{Training Objective}
\label{sec:objectives}

We retain the Scaffold-GS reconstruction and scale terms,
\begin{equation}
\begin{aligned}
\mathcal{L}_{\mathrm{rgb}}
={}&
(1-\lambda_{\mathrm{SSIM}})\,\mathcal{L}_{1}(\hat{I},I)
+\lambda_{\mathrm{SSIM}}\bigl[1-\operatorname{SSIM}(\hat{I},I)\bigr]
\\
&+\lambda_{\mathrm{s}}\,
\mathbb{E}\!\left[\prod_{d=1}^{3}s_{d}\right],
\end{aligned}
\label{eq:rgb-loss}
\end{equation}
where $s_{d}$ are scales of Gaussians retained by the opacity mask,
$\lambda_{\mathrm{SSIM}}$ is the usual trade-off between pixel $\ell_1$ and SSIM,
and $\lambda_{\mathrm{s}}$ penalizes overly large Gaussians.
An opacity penalty is applied to decoder outputs \emph{before} that mask,
\begin{equation}
\mathcal{L}_{\mathrm{op}}
=
\mathbb{E}\bigl[\max(0,\tau_{\alpha}-\alpha)\bigr],
\label{eq:opacity-loss}
\end{equation}
where $\tau_{\alpha}$ is a hinge threshold (zero in our setting, so only negative opacities are penalized).
The weight $\lambda_{\mathrm{op}}$ of this term is introduced in the total loss below.

For an internal router $\nu$, let $\mathbf{P}_{\nu}=\mathbb{E}_{i}[\mathbf{g}_{i,\nu}]$ be the mean softmax gate
and let $\mathbf{a}_{\nu}$ be the mean child-selection indicator (binary under Top-$k$, equal to $\mathbf{P}_{\nu}$ when dense).
The Switch-style load-balancing loss averages over routers of a head and then over the three heads:
\begin{equation}
\mathcal{L}_{\mathrm{load}}
=
\frac{1}{3}\sum_{h}
\frac{1}{|\mathcal{I}_{h}|}
\sum_{\nu\in\mathcal{I}_{h}}
b_{\nu}\,
\mathbf{P}_{\nu}^{\top}\mathbf{a}_{\nu},
\label{eq:load-loss}
\end{equation}
where $\mathcal{I}_{h}$ is the set of internal routers of head $h$ (a single root router in the flat MoE).
Router statistics are computed over visible anchors, unweighted by path mass.
We also penalize the entropy of the dense leaf-path distribution,
\begin{equation}
\mathcal{L}_{\mathrm{ent}}
=
-\frac{1}{3}\sum_{h}
\mathbb{E}_{i}
\left[
\sum_{\ell\in\mathcal{L}_{h}}
\tilde{\pi}_{i}^{(h)}(\ell)
\log\tilde{\pi}_{i}^{(h)}(\ell)
\right].
\label{eq:entropy-loss}
\end{equation}
Minimizing $\mathcal{L}_{\mathrm{ent}}$ encourages concentrated routing: each anchor should commit to few leaves rather than spreading mass uniformly.
The full objective is
\begin{equation}
\mathcal{L}
=
\mathcal{L}_{\mathrm{rgb}}
+\lambda_{\mathrm{op}}\mathcal{L}_{\mathrm{op}}
+\lambda_{\mathrm{load}}(t)\,\mathcal{L}_{\mathrm{load}}
+\lambda_{\mathrm{ent}}(t)\,\mathcal{L}_{\mathrm{ent}},
\label{eq:total-loss}
\end{equation}
with a curriculum on the auxiliary weights,
\begin{equation}
\lambda_{\mathrm{load}}(t)
=
\begin{cases}
\lambda_{\mathrm{load}}, & t\le T_{\mathrm{w}},\\
\rho_{\mathrm{load}}\,\lambda_{\mathrm{load}}, & t>T_{\mathrm{w}},
\end{cases}
\quad
\lambda_{\mathrm{ent}}(t)
=
\begin{cases}
0, & t\le T_{\mathrm{w}},\\
\lambda_{\mathrm{ent}}, & t>T_{\mathrm{w}}.
\end{cases}
\label{eq:loss-schedule}
\end{equation}
During warm-up, $\lambda_{\mathrm{load}}$ encourages every expert to receive some traffic, while entropy regularization is off so the router may still explore.
After warm-up, load balancing is scaled down by $\rho_{\mathrm{load}}\in(0,1)$ and entropy regularization with weight $\lambda_{\mathrm{ent}}$ is turned on, so routing may specialize.
This schedule follows the global iteration counter and is not restarted during post-growth dense windows.

\subsection{Two-Stage Training}
\label{sec:recipe}

\paragraph{Stage A: flat MoE.}
From the Scaffold-GS SfM initialization we train the $E$-expert flat MoE for $T$ iterations, including the standard anchor densification and pruning schedule.
$T$ is the length of one full training stage.

\paragraph{Stage B: growing tree MoE.}
The tree experts and root routers, together with the stored anchors (features and geometric parameters), are copied from Stage~A.
Densification and pruning are disabled.
Stage~B runs a fresh $T$-iteration schedule with its own warm-up $T_{\mathrm{w}}$.
Growth is allowed only inside a window $[T_{\mathrm{start}},T_{\mathrm{stop}}]$, and actual split attempts occur at a prescribed set of iterations $\mathcal{T}_{\mathrm{g}}$ inside that window.

\begin{algorithm}[t]
\caption{Growing Tree MoE (Stage B)}
\label{alg:tree}
\small
\begin{algorithmic}[1]
\Require Flat-MoE checkpoint; growth schedule $\mathcal{T}_{\mathrm{g}}$ and limits $(D_{\max},G_{\max},L_{\max},H_{\mathrm{evt}})$.
\Ensure Trained tree MoE.
\State Copy $E$ root leaves and the root router from Stage~A; load anchors; disable densification; reset growth statistics.
\For{$t=1,\ldots,T$}
  \State Dense routing if $t\le T_{\mathrm{w}}$ or inside a post-growth dense window; else node-wise Top-$k$.
  \State Render, compute $\mathcal{L}$, backpropagate, and record $(m_{\ell},L_{\ell},v_{\ell})$.
  \State Adam step (if $t<T$).
  \If{$t\in\mathcal{T}_{\mathrm{g}}$}
    \State Score eligible leaves with Eqs.~\eqref{eq:growth-score}--\eqref{eq:growth-gates}; relax thresholds for lagging heads.
    \State Order proposals by lag, head priority, then $S_{\ell}$; split at most $H_{\mathrm{evt}}$ heads.
    \If{any split succeeds}
      \State Reset statistics of split heads; register new Adam parameters; dense-route for $T_{\mathrm{d}}$ steps.
    \EndIf
  \EndIf
\EndFor
\State \Return trained tree MoE.
\end{algorithmic}
\end{algorithm}

\begin{algorithm}[t]
\caption{\textsc{SplitLeaf}($h,\ell$)}
\label{alg:grow}
\small
\begin{algorithmic}[1]
\State Abort if $\operatorname{depth}(\ell)\ge D_{\max}$, head $h$ has used $G_{\max}$ splits, or $|\mathcal{L}_{h}|+(G-1)>L_{\max}$.
\State Replace $\ell$ by a $G$-way linear router and $G$ children cloned from $\ell$ as in Eq.~\eqref{eq:split}.
\end{algorithmic}
\end{algorithm}

% Place in Experiments. Do not put unswept growth-score constants in the main paper.
Training length and the photometric / scale losses follow 3D-GS and Scaffold-GS.
Top-$k$ routing, noisy gating, and the Switch load-balancing form follow
Shazeer~\textit{et al.}~\cite{shazeer2017moe} and Fedus~\textit{et al.}~\cite{fedus2022switch}.
MoE width, Top-$k$, the dense warm-up, load balancing, and the split budget
are the knobs we actually vary; remaining growth-score coefficients
($\beta$, $\lambda_{v}$, $\tau_{p}$, $\tau_{R}$, $\tau_{v}$, $\tau_{S}$)
are held fixed on every scene and are not claimed as tuned contributions.

\begin{table}[t]
\centering
\caption{Settings inherited from prior work. We do not retune these per scene.}
\label{tab:hparams-inherited}
\setlength{\tabcolsep}{4pt}
\small
\begin{tabular}{lllp{0.42\linewidth}}
\toprule
Symbol & Value & Source & Role \\
\midrule
$T$ & $30\mathrm{k}$ & 3D-GS~\cite{kerbl20233d} & Iterations per stage \\
$\lambda_{\mathrm{SSIM}}$ & $0.2$ & 3D-GS~\cite{kerbl20233d} & SSIM versus $\ell_1$ in $\mathcal{L}_{\mathrm{rgb}}$ \\
$\lambda_{\mathrm{s}}$ & $0.01$ & Scaffold-GS~\cite{lu2024scaffold} & Gaussian scale regularizer \\
$G$ & $2$ & binary split & Children created by one growth event \\
$k$ (mechanism) & Top-$k$ & Switch~\cite{shazeer2017moe,fedus2022switch} & Sparse expert (or child) selection \\
$\mathcal{L}_{\mathrm{load}}$ & Switch form & Fedus~\textit{et al.}~\cite{fedus2022switch} & Load-balancing auxiliary \\
$\boldsymbol{\eta}$ & Gaussian noise & Shazeer~\textit{et al.}~\cite{shazeer2017moe} & Logit noise during dense routing \\
\bottomrule
\end{tabular}
\end{table}

\begin{table}[t]
\centering
\caption{Settings introduced in this work, kept only when a sensitivity run exists.
Mean PSNR is on the 19-scene paper split.}
\label{tab:hparams}
\setlength{\tabcolsep}{4pt}
\small
\begin{tabular}{llp{0.55\linewidth}}
\toprule
Symbol & Default & Sensitivity \\
\midrule
$E$ & $8$ & Flat MoE $E\in\{4,8\}$ with Top-$1$/$2$; default is the standard recipe, not a unique peak. \\
$k$ & $2$ & Same $E{\times}$Top-$k$ grid. Top-$2$ is the reported operating point. \\
$G_{\max}$ & $4$ & $G_{\max}\in\{0,2,4\}$. $G_{\max}{=}0$ already gives $+0.30$\,dB vs.\ flat MoE; $4$ vs.\ $0$ is $+0.005$\,dB. \\
$T_{\mathrm{w}}$ & $3\mathrm{k}$ & On vs.\ off (no curriculum). $\Delta$ vs.\ Full is $-0.01$\,dB. \\
$\lambda_{\mathrm{load}}$ & $0.005$ & On vs.\ off. $\Delta$ vs.\ Full is $0.00$\,dB; kept for Switch-style balance, not PSNR. \\
Warm-start & on & Off drops to $30.56$\,dB, i.e.\ Scaffold-GS level. \\
\bottomrule
\end{tabular}
\end{table}
